% Generated draft: requires mathematical and source review.
\documentclass[11pt]{article}
\usepackage[margin=25mm]{geometry}
\usepackage{fontspec}
\setmainfont{Latin Modern Roman}
\usepackage{amsmath,amssymb,mathtools,unicode-math}
\setmathfont{Latin Modern Math}
\usepackage{enumitem,needspace,xurl,hyperref}
\setlength{\parindent}{0pt}
\setlength{\parskip}{5pt}
\setlength{\emergencystretch}{3em}
\newcommand{\N}{\mathbb{N}}
\newcommand{\Z}{\mathbb{Z}}
\newcommand{\Q}{\mathbb{Q}}
\newcommand{\R}{\mathbb{R}}
\newcommand{\C}{\mathbb{C}}
\newcommand{\F}{\mathbb{F}}
\newcommand{\A}{\mathbb{A}}
\newcommand{\PP}{\mathbb{P}}
\newcommand{\E}{\mathbb{E}}
\newcommand{\eps}{\varepsilon}
\newcommand{\dd}{\,\mathrm{d}}
\begin{document}
\section*{UnsolvedMath ID: 30001922}
\textbf{Source problem code:} catalogue entry
\subsection{Definitions and mathematical statement}
Let \(n=\dim X\). In the usual characteristic-zero terminology, a projective minimal terminal variety is a normal projective variety with terminal singularities and nef \(\mathbb Q\)-Cartier canonical divisor \(K_X\). Write
\[
\operatorname{NS}_{\mathbb C}(X)=\operatorname{NS}(X)\otimes_{\mathbb Z}\mathbb C,
\]
where \(\operatorname{NS}(X)\) is the group of line bundles modulo algebraic equivalence. A group is virtually solvable if it has a solvable subgroup of finite index.

Let \(X\) be a projective minimal terminal variety of dimension \(n\), and let \(G\leq\operatorname{Aut}(X)\). Assume that the image of the induced representation
\[
G\longrightarrow\operatorname{GL}\bigl(\operatorname{NS}_{\mathbb C}(X)\bigr)
\]
is not virtually solvable. Write \(R(G)\) for the radical referred to in the supplied record; its precise definition is not supplied. Reading passage to finite index followed by quotient by \(R(G)\) as taking the image of a finite-index subgroup in \(G/R(G)\), assume that there exist \(G_0\leq G\) with \([G:G_0]<\infty\), an almost simple real linear algebraic group \(H\), and an injective homomorphism
\[
\iota:G_0/(G_0\cap R(G))\hookrightarrow H(\mathbb R)
\]
whose image is Zariski dense in \(H\), such that
\[
\operatorname{rank}_{\mathbb R}H=n-1.
\]
Here real rank is the dimension of a maximal \(\mathbb R\)-split torus.

\textbf{Question.} Must \(X\) be birational to a quotient of an abelian variety of dimension \(n\), without imposing \(n=3\)? Under the conventional interpretation of ``quotient'' as a finite group quotient, the requested conclusion is the existence of an abelian variety \(A\) and a finite subgroup \(F\leq\operatorname{Aut}(A)\) such that
\[
\dim A=n,\qquad X\dashrightarrow A/F\quad\text{is birational}.
\]

\textbf{Unresolved source qualifications.} The supplied formulation does not explicitly specify the ground field, define \(R(G)\), distinguish the displayed quotient from quotienting \(G_0\) by its own radical, or specify the kind of quotient of an abelian variety intended. The finite-group conclusion above is therefore conditional on that interpretation. The original question invokes the hypotheses of Theorem~3, whose text is not included. Although the registry reports that those hypotheses were restored, their completeness cannot be checked from the supplied record. These qualifications prevent treating this draft as an independently verified, fully specified formulation.
\subsection{Short English statement}
The supplied question asks whether a projective minimal terminal variety with the stated nonsolvable automorphism action and a Zariski-dense almost simple quotient of real rank \(\dim X-1\) must be birational to an abelian-variety quotient in arbitrary dimension. The ground field, radical, quotient convention, and completeness of the restored Theorem~3 hypotheses require review.
\subsection{Sources}
\begin{itemize}
\item Supplied local UnsolvedMath registry record, permanent ID \(30001922\), title ``Abelian-Quotient Models from High-Rank Group Actions,'' exact source problem code \texttt{catalogue entry}, in the Oberwolfach Reports Open Problems collection. No individual proposer is identified; the supplied proposed year is 2011.
\item Supplied source citation: \emph{Complex Algebraic Geometry (2012)}, Oberwolfach Reports, DOI \texttt{10.4171/owr/2011/47}, pp.~2669--2733. Original OWR report: \url{https://doi.org/10.4171/owr/2011/47}. The original question refers to Theorem~3 and asks for the real-rank \(n-1\) conclusion without assuming \(n=3\).
\item Supplied literature reference, labelled ``Cantat--Zeghib lattice case'': \url{https://arxiv.org/abs/1011.4768}. The record supplies no further bibliographic details or precise theorem hypotheses for this reference.
\item Literature assessment dated exactly \texttt{2026-08-21}: the supplied curation labels the question open beyond lattice cases attributed to Cantat--Zeghib and reports finding no general theorem establishing the stated birational classification. This is supplied dated curation, not an independently verified literature assessment; its status and cited claims require independent checking.
\item The registry labels the statement \texttt{corrected\_verified} and says that hypotheses suppressed by the Theorem~3 reference were restored from the official report. This is a supplied verification claim, not verification performed for this draft. No linked source was visited in preparing this statement.
\end{itemize}
\end{document}
